% ==============================================================================
% CHƯƠNG 1: MỞ ĐẦU
% ==============================================================================
\chapter{Mở đầu}
\label{chap:mo-dau}

\section{Lý do chọn đề tài}
\label{sec:ly-do}

Trình bày bối cảnh và vấn đề đặt ra. Nêu \nhanmanh{khoảng trống} hoặc
\nhanmanh{nhu cầu thực tế} khiến đề tài đáng thực hiện, có dẫn chứng từ tài liệu
\cite{acemoglu2022, brynjolfsson2025}. Khi nhắc một khái niệm lần đầu, dùng
\thuatngu{trí tuệ nhân tạo}[Artificial Intelligence] để làm nổi bật thuật ngữ.

Đoạn thứ hai giải thích vì sao vấn đề quan trọng với người đọc (sinh viên, doanh
nghiệp, cơ quan quản lý\dots).

\section{Mục tiêu báo cáo}
\label{sec:muc-tieu}

Báo cáo hướng tới các mục tiêu sau:
\begin{itemize}
  \item Mục tiêu 1: phân tích \dots
  \item Mục tiêu 2: trình bày kết quả định lượng về \dots
  \item Mục tiêu 3: đề xuất giải pháp cho \dots
\end{itemize}

\section{Phạm vi và phương pháp}
\label{sec:pham-vi}

Nêu rõ đối tượng, phạm vi (thời gian, không gian, dữ liệu) và phương pháp chính.

\begin{luuy}[Giới hạn phạm vi]
  Ghi rõ những gì báo cáo \quantrong{không} đề cập để người đọc không hiểu sai
  phạm vi kết luận.
\end{luuy}

\section{Bố cục báo cáo}
\label{sec:bo-cuc}

Ngoài phần mở đầu, báo cáo gồm các chương sau:
\begin{itemize}
  \item \Chuong{chap:co-so}: cơ sở lý thuyết và các khái niệm nền tảng.
  \item \Chuong{chap:phuong-phap}: dữ liệu và phương pháp thực hiện.
  \item \Chuong{chap:ket-qua}: kết quả và thảo luận.
  \item \Chuong{chap:ket-luan}: kết luận và kiến nghị.
\end{itemize}
